\section*{第 \ref{ch:g11:mechanical-energy} 章——机械能及其守恒}

\begin{solution}{exo:g11:mechanical-energy:1}
(a) $v = \qty{6.94}{m/s}$：
$E_k = \tfrac12 \times 90 \times 6.94^2 \approx \qty{2.2e3}{J}$。
(b) $v = \qty{36.1}{m/s}$：
$E_k = \tfrac12 \times 1200 \times 36.1^2 \approx \qty{7.8e5}{J}$。
(c) $\times 4$；$\times 9$——平方关系。
\end{solution}

\begin{solution}{exo:g11:mechanical-energy:2}
$\Delta E_p = 75 \times 9.81 \times 850 \approx \qty{6.3e5}{J}$。不对：
参考平面的选取使两个值同时平移相同的量。不对：只有高度
差参与计算。
\end{solution}

\begin{solution}{exo:g11:mechanical-energy:3}
$v = \sqrt{2 \times 9.81 \times 20} = \qty{19.8}{m/s} \approx
\qty{71}{km/h}$。$mgh = \tfrac12 mv^2$，质量消去。
\end{solution}

\begin{solution}{exo:g11:mechanical-energy:4}
$W = \Delta E_k = 0 - \tfrac12 \times 1200 \times 25^2 =
\qty{-3.75e5}{J}$；$F = \num{3.75e5}/50 = \qty{7.5e3}{N}$。
\end{solution}

\begin{solution}{exo:g11:mechanical-energy:5}
$v = \sqrt{2 \times 9.81 \times 0.080} \approx \qty{1.3}{m/s}$；它会
回升至 \qty{8.0}{cm}（机械能守恒）。
\end{solution}

\begin{solution}{exo:g11:mechanical-energy:6}
$\num{1.5e9}/\num{2.6e3} \approx \num{5.9e5}$：约六十
万颗子弹。$h = 88.9^2/(2 \times 9.81) \approx \qty{4.0e2}{m}$
——其自身速度足以将列车抬升 $400$ 米。
\end{solution}

\begin{solution}{exo:g11:mechanical-energy:7}
速率相同，$v = \sqrt{2 \times 9.81 \times 3.0} \approx \qty{7.7}{m/s}$
（只与下降高度有关）。陡坡滑道用时更短：它能更快达到较高速率，
且路程更短。
\end{solution}

\begin{solution}{exo:g11:mechanical-energy:8}
$h = v^2/2g = 144/19.62 \approx \qty{7.3}{m}$。在 \qty{4.0}{m} 处：
$v = \sqrt{144 - 2 \times 9.81 \times 4.0} \approx \qty{8.1}{m/s}$ ---
两个方向上结果相同：$E_m$ 只取决于高度，与运动方向无关。
\end{solution}

\begin{solution}{exo:g11:mechanical-energy:9}
$v = \sqrt{23^2 + 2 \times 9.81 \times 35} = \sqrt{1216} \approx
\qty{35}{m/s}$。实际值更小：空气阻力耗散了部分
$E_m$。
\end{solution}

\begin{solution}{exo:g11:mechanical-energy:10}
$\Delta E_m = \tfrac12 \times 45 \times 8.0^2 - 45 \times 9.81 \times
6.0 = 1440 - 2649 \approx \qty{-1.2e3}{J}$；
$f = 1209/15 \approx \qty{81}{N}$。转化为热量，进入滑雪板与
雪面之中。
\end{solution}

\begin{solution}{exo:g11:mechanical-energy:11}
$v = \sqrt{2 \times 9.81 \times 0.15} \approx \qty{1.7}{m/s}$；它会
再次升至 \qty{15}{cm}。钉子改变了路径（圆周更小），而
不改变能量：高度与速率均不受影响。
\end{solution}

\begin{solution}{exo:g11:mechanical-energy:12}
$v = \sqrt{2 \times 9.81 \times 2.5} = \qty{7.0}{m/s}$。上升高度与
$E_m$ 成正比，因此每次往返 $\times 0.9$：$2.25$、$2.03$、$1.82$、
$1.64$、$1.48$、$1.33$、$1.20$、$1.08$、$0.97$——在第
$9$ 次往返时首次低于 \qty{1.0}{m}。
\end{solution}

\begin{solution}{exo:g11:mechanical-energy:13}
初始 $E_m = \tfrac12 m \times 15^2 + m \times 9.81 \times 45$ 相同，无论角度如何，
因此落地速率为
$v = \sqrt{225 + 2 \times 9.81 \times 45} = \sqrt{1108} \approx
\qty{33}{m/s}$。角度决定落地速度的\emph{方向}、飞行时间与
射程——不决定落地速率。
\end{solution}

\begin{solution}{exo:g11:mechanical-energy:14}
$\tfrac12 mv^2 = mgh - f\ell = 80 \times 9.81 \times 120 - 65 \times
800 = \num{94176} - \num{52000} = \qty{4.22e4}{J}$，于是
$v = \sqrt{2 \times \num{42176}/80} \approx \qty{32}{m/s}$
（\qty{117}{km/h}）。耗散比例：$52000/94176 \approx 55\%$。
\end{solution}

\begin{solution}{exo:g11:mechanical-energy:15}
$m = \qty{2.0e9}{kg}$：
$E_p = \num{2.0e9} \times 9.81 \times 300 \approx \qty{5.9e12}{J}
= \num{5.9e12}/\num{3.6e6} \approx \qty{1.6e6}{kW.h}$。回收的能量：
$0.85 \times \num{1.64e6} \approx \qty{1.4e6}{kW.h}$——约可供
$\num{1.4e5}$ 户家庭用电一天。
\end{solution}

\begin{solution}{pb:g11:mechanical-energy:1}
\textbf{1.} $v_{\text{顶}} = \sqrt{gR} = \sqrt{9.81 \times 8.0}
\approx \qty{8.9}{m/s}$。

\textbf{2.} $E_m = mg\,(2R) + \tfrac12 m gR = \num{3.14e5} +
\num{7.8e4} \approx \qty{3.92e5}{J}$。

\textbf{3.} $mgH_{\min} = E_m$ 得 $H_{\min} = \qty{20}{m}$；
代数上 $H_{\min} = 2R + \frac{gR}{2g} = \frac{5R}{2}$。

\textbf{4.} $v_{\text{顶}} = \sqrt{2g(H - 2R)} =
\sqrt{2 \times 9.81 \times 8.0} \approx \qty{12.5}{m/s}$；
$v_{\text{顶}}^2/(gR) = 2.0$——是最小值的两倍。

\textbf{5.} 各能量均与 $m$ 成正比，故 $m$ 约去：
速率相同，安全性亦相同，无论满载或空载。

\textbf{6.} $v = \sqrt{2 \times 9.81 \times 24} \approx \qty{21.7}{m/s}
\approx \qty{78}{km/h}$.

\textbf{7.} $v = \sqrt{2 \times 9.81 \times (24 - 8.0)} \approx
\qty{17.7}{m/s}$.

\textbf{8.} $mgH = \qty{4.71e5}{J}$。$(E_p, E_k)$：出发点
$(\num{4.71e5}, 0)$；底部 $(0, \num{4.71e5})$；侧面
$(\num{1.57e5}, \num{3.14e5})$；顶部 $(\num{3.14e5}, \num{1.57e5})$——
每组两数之和均为 \qty{4.71e5}{J}。

\textbf{9.} $v = \sqrt{2 \times 9.81 \times (24 - 12)} \approx
\qty{15.3}{m/s}$.

\textbf{10.} $v(z) = \sqrt{2g(H - z)}$：速率在轨道的
最低点最大，环路顶部最小。

\textbf{11.} 高度相同，故 $\Delta E_p = 0$，且
$\Delta E_m = \Delta E_k = \tfrac12 \times 2000 \times (19.2^2 -
20.4^2) \approx \qty{-4.75e4}{J}$：此处的安排使亏损纯为
动能，两个速度计即可读出。

\textbf{12.} $W(\vect f) = \qty{-4.75e4}{J}$;
$f = \num{4.75e4}/60 \approx \qty{7.9e2}{N}$.

\textbf{13.} 热能，存在于车轮、轨道和空气中。

\textbf{14.} 在第一个传感器处，$E_m$ 为 $\num{4.16e5} + \num{3.9e4} =
\qty{4.55e5}{J}$；损失 $\num{4.75e4}/\num{4.55e5} \approx 10\%$，
每 \qty{60}{m}。在 $H_{\min}$ 时，环路顶部\emph{恰好}
达到 $\sqrt{gR}$ 的速率；任何损失都会使其降至安全速率以下。

\textbf{15.} 每 \qty{60}{m} 损失 $\approx 10\%$，约 $15\%$ 的
$mgH$（$\approx \qty{7.1e4}{J}$）已损失，因此 $E_m \approx \qty{4.00e5}{J}$，
仍高于所需的 \qty{3.92e5}{J}——安全，仅余 $2\%$
的余量。这一余量并非奢侈。

\textbf{16.} $E_k = \tfrac12 \times 2000 \times 18.0^2 =
\qty{3.24e5}{J}$.

\textbf{17.} $F = \num{3.24e5}/45 = \qty{7.2e3}{N}$.

\textbf{18.} $a = 7200/2000 = \qty{3.6}{m/s^2} \approx 0.37\,g$:
感觉较强但仍然舒适。

\textbf{19.} 轨道摩擦损耗：$mgH - \num{3.24e5} = \num{470880} -
\num{324000} = \qty{1.47e5}{J}$。核对：$\num{1.47e5} + \num{3.24e5} =
\num{4.71e5} = mgH$——每一焦耳均已计入。

\textbf{20.} 出发高度 $H = \qty{24}{m}$ 决定了环路的
安全性、任一点的速率 $\sqrt{2g(H - z)}$、摩擦损耗与
制动力；忽略摩擦的代价，便是余量——即 $\qty{4}{m}$
的山坡高度，被有摩擦的现实世界悄然耗去。
\end{solution}
